\begin{Exercise}[label={ex:massimo}]
  Scrivere un programma che chieda all'utente di inserire un numero variabile (al massimo $10$) di valori reali. Il programma dovr\`a assicurarsi che l'utente non possa inserire pi\`u di $10$ valori.

  Il programma dovr\`a cercare, durante l'acquisizione, la posizione del minimo e del massimo.

  Infine, il programma dovr\`a stampare un sommario dei valori inseriti, con l'indicazione della posizione del minimo (indicata con un segno $-$) e del massimo (indicata con un segno $+$), il tutto formattato come indicato di seguito.

\begin{lstlisting}[style=console]
Quanti elementi vuoi inserire (max: 10)? 100
Il valore inserito è troppo grande!
Quanti elementi vuoi inserire (max: 10)? 4
Inserire il 1o valore: 100
Inserire il 2o valore: 20
Inserire il 3o valore: -1
Inserire il 4o valore: 5.3
  +| 100.00
   | 20.00
-  | -1.00
   | 5.30
\end{lstlisting}
\end{Exercise}

\begin{Answer}[ref={ex:massimo}]
  \lstinputlisting[style=cstylesmall]{src/massimo.c}
\end{Answer}
